\usepackage{xeCJK}
\usepackage{geometry}
\geometry{left=2cm,right=2cm,top=2.5cm,bottom=2.5cm}
\setlength{\headheight}{12.7pt}

\usepackage{titlesec}
\titleformat{\section}
  {\normalfont\large\bfseries}
  {\thesection}
  {1em}
  {}
\titlespacing*{\section}
  {0pt}
  {3.5ex plus 1ex minus .2ex}
  {2.3ex plus .2ex}

\usepackage{amsmath}
\usepackage{amssymb}
\usepackage{amsthm}
\usepackage{enumitem}
\setlist[enumerate]{itemsep=0pt,topsep=0pt,parsep=0pt,leftmargin=0pt,itemindent=2.5em}
\setlist[itemize]{itemsep=0pt,topsep=0pt,parsep=0pt,leftmargin=0pt,itemindent=2.5em}
\usepackage{fancyhdr}
\pagestyle{fancy}
\fancyhf{}
\usepackage{graphicx}
\usepackage{hyperref}
\usepackage{indentfirst}
\usepackage{lmodern}


\begin{document}
\setlength{\abovedisplayskip}{1pt}
\setlength{\belowdisplayskip}{1pt}

\section{一阶语言的公式}

\hypertarget{sec:定义1.1}{}
\noindent\textbf{定义1.1 (一阶语言的公式)}

给定\href{01 一阶语言#nameddest=sec:定义1.1}{一阶语言} $\mathcal{L}=
(\mathcal{V},\mathcal{C},(n_f)_{f\in\mathcal{F}},(n_r)_{r\in\mathcal{R}})$,
在$\Sigma_{\mathcal{L}}^*$中定义:
\begin{enumerate}
    \item 原子公式集
    \[
        F_0=\{t_1\equiv t_2\mid t_1,t_2\in T^{\mathcal{L}}\}\cup
        \{rt_1\cdots t_n\mid n\in\mathbb{N}^+,
        r\in\mathcal{R},n_r=n,t_1,\cdots,t_n\in T^{\mathcal{L}}\},
    \]
    \item 归纳规则:
    \[\begin{aligned}
        &\mathrm{Neg}:&&\Sigma_{\mathcal{L}}^*\to\Sigma_{\mathcal{L}}^*,
        \quad&&\varphi_1\mapsto\neg\varphi_1,\\[-3pt]
        &\mathrm{To}:&&(\Sigma_{\mathcal{L}}^*)^2\to\Sigma_{\mathcal{L}}^*,
        \quad&&(\varphi_1,\varphi_2)\mapsto(\varphi_1\to\varphi_2),\\[-3pt]
        &\mathrm{Exists}^x:&&\Sigma_{\mathcal{L}}^*\to\Sigma_{\mathcal{L}}^*,
        \quad&&\varphi_1\mapsto\exists x\varphi_1
        \,(\text{对每个变元}\,x\in\mathcal{V}).
    \end{aligned}\]
\end{enumerate}
从$F_0$和归纳规则可以\href{01 一阶语言#nameddest=sec:定义A1.1}{归纳定义}\textbf{公式集}
$F^{\mathcal{L}}$, 称其中的字符串为$\mathcal{L}$的\textbf{公式}.

\vspace{10pt}

使用\href{01 一阶语言#nameddest=sec:定理A1.2}{归纳证明}显然有以下的定理.
\begin{itemize}
    \item 公式不是空字符串, 并且含有$\mathcal{R}$或$\mathrm{Log}$中的符号;
    \item 如果公式$\varphi$的首字符是常元、变元或函数符号, 则存在项$t_1,t_2$使得
    $\varphi=t_1\equiv t_2$;
    \item 如果公式$\varphi$的首字符是$n$元关系符号$r$, 则存在项$t_1,\cdots,t_n$使得
    $\varphi=rt_1\cdots t_n$;
    \item 如果公式$\varphi$的首字符是$\neg$, 则存在公式$\varphi_1$使得
    $\varphi=\neg\varphi_1$;
    \item 如果公式$\varphi$的首字符是$($, 则存在公式$\varphi_1,\varphi_2$使得
    $\varphi=(\varphi_1\to\varphi_2)$;
    \item 如果公式$\varphi$的首字符是$\exists$, 则存在变元$x$和公式$\varphi_1$使得
    $\varphi=\exists x\varphi_1$.
\end{itemize}

\vspace{10pt}

\hypertarget{sec:定理1.1}{}
\noindent\textbf{定理1.1}

如果两个公式有前缀关系, 那么它们相同.

\noindent{\kaishu 证明.}

定义命题$P(\varphi)$:
对任意的公式$\psi$, 如果$\varphi\sim\psi$, 那么$\varphi=\psi$.

接下来归纳证明即可.

\begin{enumerate}
    \item 如果存在项$t_1,t_2$使得$\varphi=t_1\equiv t_2$,
    假设$\varphi\sim\psi$, 那么存在项$s_1,s_2$使得$\psi=s_1\equiv s_2$.

    \hspace{2.17em}
    所以$t_1=s_1$, 进一步$t_2=s_2$. 即$\varphi=\psi$.

    \item 如果存在$n$元关系符号$r$和项$t_1,\cdots,t_n$使得$\varphi=rt_1\cdots t_n$,
    假设$\varphi\sim\psi$, 那么存在项$s_1,\cdots,s_n$使得$\psi=rs_1\cdots s_n$.

    \hspace{2.17em}
    所以各个$t_i=s_i$, 即$\varphi=\psi$.

    \item 假设$\varphi_1$满足命题$P$, 下证$\varphi=\neg\varphi_1$也满足命题$P$.
    
    \hspace{2.17em}
    如果$\varphi\sim\psi$, 那么存在公式$\psi_1$使得$\psi=\neg\psi_1$.

    \hspace{2.17em}
    所以$\varphi_1\sim\psi_1$. 依假设有$\varphi_1=\psi_1$, 于是$\varphi=\psi$.

    \item 假设$\varphi_1,\varphi_2$满足命题$P$,
    下证$\varphi=(\varphi_1\to\varphi_2)$也满足命题$P$.

    \hspace{2.17em}
    如果$\varphi\sim\psi$,
    那么存在公式$\psi_1,\psi_2$使得$\psi=(\psi_1\to\psi_2)$.

    \hspace{2.17em}
    所以$\varphi_1\sim\psi_1$, 依假设有$\varphi_1=\psi_1$.
    进一步$\varphi_2\sim\psi_2$, 依假设有$\varphi_2=\psi_2$. 于是$\varphi=\psi$.

    \item 对任意的变元$x$, 假设$\varphi_1$满足命题$P$,
    下证$\varphi=\exists x\varphi_1$也满足命题$P$.

    \hspace{2.17em}
    如果$\varphi\sim\psi$,
    那么存在变元$y$和公式$\psi_1$使得$\psi=\exists y\psi_1$.

    \hspace{2.17em}
    所以$x=y$, 进一步$\varphi_1\sim\psi_1$, 依假设有$\varphi_1=\psi_1$.
    于是$\varphi=\psi$. \hfill $\qedsymbol$
\end{enumerate}





\newpage

\hypertarget{sec:定理1.2}{}
\noindent\textbf{定理1.2 (公式的唯一可读性)}

对任意的公式$\varphi$, 下面五条有且只有一条成立:
\begin{enumerate}
    \item 存在唯一的项$t_1,t_2$使得$\varphi=t_1\equiv t_2$;
    \item 存在唯一的关系符号$r$和项$t_1,\cdots,t_n$使得$\varphi=rt_1\cdots t_n$;
    \item 存在唯一的公式$\varphi_1$使得$\varphi=\neg\varphi_1$;
    \item 存在唯一的公式$\varphi_1,\varphi_2$使得$\varphi=(\varphi_1\to\varphi_2)$;
    \item 存在唯一的变元$x$和公式$\varphi_1$使得$\varphi=\exists x\varphi_1$.
\end{enumerate}

\vspace{10pt}

\hypertarget{sec:定理1.3}{}
\noindent\textbf{定理1.3 (公式的映射)}

给定集合$D$, 以及:
\begin{enumerate}
    \item 对任意的原子公式$\varphi$指定一个元素$\mathrm{i}(\varphi)\in D$,
    \item 指定一个映射$\mathrm{i}(\neg):D\to D$,
    \item 指定一个映射$\mathrm{i}(\to):D^2\to D$,
    \item 对任意的变元$x$指定一个映射$\mathrm{i}(x):D\to D$,
\end{enumerate}
那么存在唯一的映射$I:F^{\mathcal{L}}\to D$使得
\begin{enumerate}
    \item 对任意的原子公式$\varphi$有$I(\varphi)=\mathrm{i}(\varphi)$,
    \item 对任意的公式$\varphi$有$I(\neg\varphi)=\mathrm{i}(\neg)(I(\varphi))$,
    \item 对任意的公式$\varphi,\psi$有
    $I((\varphi\to\psi))=\mathrm{i}(\to)(I(\varphi),I(\psi))$,
    \item 对任意的变元$x$和公式$\varphi$有
    $I(\exists x\varphi)=\mathrm{i}(x)(I(\varphi))$.
\end{enumerate}

\noindent{\kaishu 证明.}

在集合$F^{\mathcal{L}}\times D$中进行归纳定义.
\begin{enumerate}
    \item 初始集$I_0=\{(\varphi,\mathrm{i}(\varphi))\mid \varphi\in F_0\}$,
    \item 归纳规则:
    \[\begin{aligned}
        &\mathrm{Neg}_{\mathrm{i}}:&&
        F^{\mathcal{L}}\times D\to F^{\mathcal{L}}\times D,\quad&&
        (\varphi,d)\mapsto(\neg\varphi,\mathrm{i}(\neg)(d)),\\[-3pt]
        &\mathrm{To}_{\mathrm{i}}:&&
        (F^{\mathcal{L}}\times D)^2\to F^{\mathcal{L}}\times D,\quad&&
        ((\varphi,d),(\psi,d'))\mapsto((\varphi\to\psi),\mathrm{i}(\to)(d,d')),\\[-3pt]
        &\mathrm{Exists}_{\mathrm{i}}^x:&&
        F^{\mathcal{L}}\times D\to F^{\mathcal{L}}\times D,\quad&&
        (\varphi,d)\mapsto(\exists x\varphi,\mathrm{i}(x)(d))
        \,(\text{对每个变元}\,x\in\mathcal{V}).
    \end{aligned}\]
\end{enumerate}
从$I_0$和归纳规则可以归纳定义$I\subset F^{\mathcal{L}}\times D$.
下面证明$I$就是要找的映射.

定义命题$P(\varphi)$: $\varphi$在$I$中存在唯一的像. 归纳证明即可.

\begin{enumerate}
    \item 如果$\varphi$是原子公式, 显然它有唯一的像$\mathrm{i}(\varphi)$.
    
    \item 假设$\varphi_1$满足命题$P$, 下证$\varphi=\neg\varphi_1$也满足命题$P$.
    
    \hspace{2.17em}
    记$\varphi_1$的唯一像为$d$, 那么易证$\varphi$有唯一的像$\mathrm{i}(\neg)(d)$.

    \item 假设$\varphi_1,\varphi_2$满足命题$P$,
    下证$\varphi=(\varphi_1\to\varphi_2)$也满足命题$P$.

    \hspace{2.17em}
    记$\varphi_1,\varphi_2$的唯一像分别为$d,d'$,
    那么易证$\varphi$有唯一的像$\mathrm{i}(\to)(d,d')$.

    \item 对任意的变元$x$, 假设$\varphi_1$满足命题$P$,
    下证$\varphi=\exists x\varphi_1$也满足命题$P$.

    \hspace{2.17em}
    记$\varphi_1$的唯一像为$d$, 那么易证$\varphi$存在唯一的像$\mathrm{i}(x)(d)$.
\end{enumerate}

所以$I:F^{\mathcal{L}}\to D$, 易证它满足条件. 唯一性易证. \hfill $\qedsymbol$





\newpage
\section{自由变元}

\hypertarget{sec:定义2.1}{}
\noindent\textbf{定义2.1 (自由变元)}

给定一阶语言$\mathcal{L}$, 在$\mathcal{P}(\mathcal{V})$中定义公式的映射$\mathrm{FV}$:
\begin{enumerate}
    \item 对任意的项$t_1,t_2$,
    $\mathrm{FV}(t_1\equiv t_2)=\mathrm{var}(t_1)\cup\mathrm{var}(t_2)$,
    \item 对任意的$n$元关系符号$r$和项$t_1,\cdots,t_n$,
    $\mathrm{FV}(rt_1\cdots t_n)=\mathrm{var}(t_1)\cup\cdots\cup\mathrm{var}(t_n)$,
    \item 对任意的公式$\varphi$,
    $\mathrm{FV}(\neg\varphi)=\mathrm{FV}(\varphi)$,
    \item 对任意的公式$\varphi,\psi$,
    $\mathrm{FV}((\varphi\to\psi))=\mathrm{FV}(\varphi)\cup\mathrm{FV}(\psi)$,
    \item 对任意的变元$x$和公式$\varphi$,
    $\mathrm{FV}(\exists x\varphi)=\mathrm{FV}(\varphi)\setminus\{x\}$.
\end{enumerate}
称$\mathrm{FV}(\varphi)$为$\varphi$的\textbf{自由变元集},
称其中的元素为$\varphi$的自由变元.

\vspace{10pt}

\hypertarget{sec:定义2.2}{}
\noindent\textbf{定义2.2 (约束出现)}

给定一阶语言$\mathcal{L}$, 在$\mathcal{P}(\mathcal{V})$中定义公式的映射$\mathrm{bnd}$:
\begin{enumerate}
    \item 对任意的项$t_1,t_2$,
    $\mathrm{bnd}(t_1\equiv t_2)=\varnothing$,
    \item 对任意的$n$元关系符号$r$和项$t_1,\cdots,t_n$,
    $\mathrm{bnd}(rt_1\cdots t_n)=\varnothing$,
    \item 对任意的公式$\varphi$,
    $\mathrm{bnd}(\neg\varphi)=\mathrm{bnd}(\varphi)$,
    \item 对任意的公式$\varphi,\psi$,
    $\mathrm{bnd}((\varphi\to\psi))=\mathrm{bnd}(\varphi)\cup\mathrm{bnd}(\psi)$,
    \item 对任意的变元$x$和公式$\varphi$,
    $\mathrm{bnd}(\exists x\varphi)=\mathrm{bnd}(\varphi)\cup\{x\}$.
\end{enumerate}
称$\mathrm{bnd}(\varphi)$中的元素是在$\varphi$中约束出现的变元.

\vspace{10pt}

类似地可以定义公式$\varphi$的常元集$\mathrm{const}(\varphi)$,
记$\mathrm{at}=\mathrm{FV}\cup\mathrm{const}$. 易证它们都是有限集.

\vspace{10pt}

\hypertarget{sec:定义2.3}{}
\noindent\textbf{定义2.3 (语句)}

对任意的公式$\varphi$, 如果$\mathrm{FV}(\varphi)=\varnothing$,
即$\varphi$中不含自由变元, 就称$\varphi$是\textbf{语句}.

\vspace{20pt}





\section{扩展逻辑符号}

\hypertarget{sec:定义3.1}{}
\noindent\textbf{定义3.1 (扩展逻辑符号)}

为了表达方便, 常常使用$\land,\lor,\leftrightarrow,\forall$这些符号来生成公式.
具体的规则如下.
\begin{enumerate}
    \item $(\varphi\land\psi)\overset{\mathrm{def}}{=}\neg(\varphi\to\neg\psi)$,
    \item $(\varphi\lor\psi)\overset{\mathrm{def}}{=}(\neg\varphi\to\psi)$,
    \item $(\varphi\leftrightarrow\psi)\overset{\mathrm{def}}{=}
    ((\varphi\to\psi)\land(\psi\to\varphi))$,
    \item $\forall x\varphi\overset{\mathrm{def}}{=}\neg\exists x\neg\varphi$.
\end{enumerate}

\vspace{10pt}

利用这些符号可以方便地表达全称公式的定义.

\end{document}